% =====================================================================
% GreenFieldData PhD Position L – Research Proposal
% Candidate: Partha Pratim Saha
% Title: Conversational AI for Sustainable Agriculture:
%        Natural Language Interfaces over Robotic & Analytical Farming Systems
% Submitted: 11 May 2026 | Citation style: IEEE
% =====================================================================
\documentclass[11pt,a4paper]{article}

\usepackage[margin=2.0cm]{geometry}
\usepackage{times}
\usepackage{microtype}
\usepackage{setspace}
\usepackage{titlesec}
\usepackage{enumitem}
\usepackage{graphicx}
\usepackage{booktabs}
\usepackage{array}
\usepackage{xcolor}
\usepackage{hyperref}
\usepackage{cite}
\usepackage{amsmath}
\usepackage{pgfgantt}
\usepackage{rotating}
\usepackage{multirow}
\usepackage{caption}
\usepackage{subcaption}
\usepackage{fancyhdr}
\usepackage{mdframed}
\usepackage{tabularx}
\usepackage[T1]{fontenc}

\hypersetup{
  colorlinks=true, linkcolor=blue!70!black,
  citecolor=blue!60!black, urlcolor=teal
}

\definecolor{agrigreen}{RGB}{34,139,34}
\definecolor{techblue}{RGB}{0,70,127}
\definecolor{warnorange}{RGB}{210,105,30}

\titleformat{\section}{\large\bfseries\color{techblue}}{}{0em}{}[\titlerule]
\titleformat{\subsection}{\normalsize\bfseries\color{agrigreen}}{}{0em}{}
\titleformat{\subsubsection}{\normalsize\itshape}{}{0em}{}

\setlength{\parskip}{4pt}
\setlength{\parindent}{0pt}
\onehalfspacing

\pagestyle{fancy}
\fancyhf{}
\lhead{\small\textit{GreenFieldData PhD-L | Partha Pratim Saha}}
\rhead{\small\textit{Conversational AI for Sustainable Agriculture}}
\cfoot{\small\thepage}

% =====================================================================
\begin{document}

% ---- Title Block ----------------------------------------------------
\begin{center}
  {\LARGE\bfseries\color{techblue}
    Conversational AI for Sustainable Agriculture:\\[4pt]
    Natural Language Interfaces over Robotic and Analytical Farming Systems}\\[10pt]
  {\large Partha Pratim Saha}\\[2pt]
  {\normalsize
    \texttt{technical.partha@gmail.com}
    \quad|\quad
    GreenFieldData MSCA Joint Doctorate — Position L\\[2pt]
    Supervisors: Dr.\ Genoveva Vargas-Solar (CNRS/UCBLyon1) \& Prof.\ Roberto Oberti (UniMI)}\\[6pt]
  {\small Submitted: 11 May 2026}
\end{center}
\hrule height 1.2pt
\vspace{6pt}

% ---- Abstract Box ---------------------------------------------------
\begin{mdframed}[backgroundcolor=gray!8, linewidth=0.5pt]
\textbf{Abstract.}
Agricultural robotics and IoT generate heterogeneous, high-velocity data streams that require expert
programmers to interpret — effectively locking non-expert farmers and agronomists outside the
decision loop. This proposal addresses that gap with \textbf{AgriTalk}: a Metaflow-orchestrated,
LLM-powered conversational layer that translates natural language into calibrated robotic commands
and analytical queries across field robots, IoT sensors, and streaming analytics infrastructure.
The focal scenario is \emph{precision pesticide spraying} — a safety-critical, latency-sensitive task that
demands semantic grounding, uncertainty-aware human confirmation, and auditable decision trails.
The three-year plan spans NLP pipeline engineering at UCBLyon1 and robotic integration at UniMI,
yielding reproducible, responsible, and publication-ready research contributions.
\end{mdframed}
\vspace{4pt}

% =====================================================================
\section{1.\ Research Framing: Problem, Motivation, and Scientific Gap}
% =====================================================================

\subsection{1.1 Why Conversational AI for Agriculture?}

The global agri-robot market grew at 19.3\% CAGR from 2021 to 2022, reaching USD 5.9B, and is
forecast to reach USD 11.9B by 2026~\cite{phd_adv_2026}. Yet this growth creates a paradox:
the most sophisticated precision tools — GPS-guided sprayers, aerial drones, multi-modal sensing
platforms — remain inaccessible to the very farmers they are meant to serve. Current interfaces
presuppose Python programming fluency, low-level API literacy, and systems administration skills
comparable to professional engineers~\cite{phd_adv_2026}.

Natural Language Processing (NLP) and Large Language Models (LLMs) offer a structurally
compelling solution. A farmer should be able to ask \emph{``How much pesticide was sprayed on the
northern field yesterday, and under what soil-humidity conditions?''} and receive a grounded, traceable
answer — without writing a single line of SQL or YAML. Conversational AI has demonstrated this
pattern in information retrieval~\cite{wagner_2025}, healthcare assistants~\cite{brown_gpt3_2020},
and customer-facing APIs~\cite{chen_2021}. Agriculture, with its heterogeneous sensing modalities,
physical safety constraints, and non-expert user population, poses a richer and under-explored frontier.

\subsection{1.2 The Scientific Gap}

Three structural problems remain unsolved:
\begin{enumerate}[leftmargin=1.2em,itemsep=1pt]
  \item \textbf{Semantic grounding under domain uncertainty.} LLMs trained on internet text lack
  grounded knowledge of crop phenology, robot kinematic constraints, or
  pesticide biochemistry. Intent recognition must be coupled with structured domain
  knowledge bases and uncertainty quantification.
  \item \textbf{Latency-aware edge–cloud deployment.} Many agricultural decisions (e.g., abort a
  spray run on wind-speed spike) require sub-second response at the edge, not cloud
  round-trip latency. Existing conversational AI systems are cloud-centric.
  \item \textbf{Responsible automation.} LLM outputs in safety-critical systems can be
  confidently wrong — \emph{hallucinating} pesticide dosages or missing a no-spray zone.
  Human-in-the-loop (HITL) confirmation, role-based permissions, and uncertainty
  reporting are therefore non-negotiable design requirements~\cite{amershi_2019,kaur_2020}.
\end{enumerate}

\subsection{1.3 Focus Scenario: Precision Pesticide Spraying}

I focus on precision pesticide spraying as the primary research scenario because it is:
(a)~safety-critical (wrong dosage affects human health and ecosystems);
(b)~data-rich (IoT sensors for humidity, wind, crop health; drone imagery; robot GPS trails);
(c)~latency-sensitive (decisions must reach actuators within 200ms at edge);
(d)~explicitly identified in the GreenFieldData position as a high-priority use case~\cite{phd_adv_2026};
and (e)~directly connected to Prof.\ Oberti's expertise in precision crop protection systems~\cite{oberti_2025}.

This conversational system — which I call \textbf{AgriTalk} — does not replace agronomist judgment.
It acts as a transparent, auditable \emph{interface} that amplifies human decision-making by
surfacing relevant data, quantifying uncertainty, and requesting explicit human confirmation before
actuating any irreversible robotic action.

% =====================================================================
\section{2.\ Proposed PhD Tasks: Five Research and Engineering Milestones}
% =====================================================================

\begin{mdframed}[backgroundcolor=blue!4, linewidth=0.4pt]
\textbf{Design philosophy.} Each task is a self-contained contribution with its own evaluation
protocol, while feeding into the next task as infrastructure. This mirrors the Metaflow DAG
philosophy: steps are modular, versioned, and independently reproducible~\cite{tagliabue_metaflow_2023}.
\end{mdframed}

% ---- Task 1 ---------------------------------------------------------
\subsection{Task 1 — Heterogeneous IoT Stream Integration \& Semantic Schema Registry
  \textnormal{[Months 1--8, UCBLyon1]}}

\textbf{Objective.}
Build a unified data ingestion and schema registry layer that ingests heterogeneous field streams
(soil sensors, drone imagery, robotic telemetry, weather APIs) and exposes them as queryable,
semantically annotated artifacts.

\textbf{Input data.}
Simulated PANGAEA-format soil datasets~\cite{kamilaris_2017}; ACRE benchmark drone imagery;
FieldRobot Event datasets; OpenWeather API.

\textbf{Expected output.}
A Metaflow \texttt{@step}-decorated ingestion DAG that: (a)~reads raw streams via Apache Kafka
and Spark Structured Streaming; (b)~maps each field to a semantic schema (OWL ontology for
agriculture — AGROVOC); (c)~persists versioned data artifacts in S3/MinIO.

\textbf{Relation to overall objective.}
Without a clean, semantically labeled data layer, NLP grounding in Task 2 is impossible.
This task directly addresses the streaming data challenge stated in~\cite{phd_adv_2026} and
benchmarks data pipeline quality using Great Expectations~\cite{tagliabue_metaflow_2023}.

% ---- Task 2 ---------------------------------------------------------
\subsection{Task 2 — Intent Recognition, Semantic Grounding \& NL-to-Command Compilation
  \textnormal{[Months 5--14, UCBLyon1]}}

\textbf{Objective.}
Design a three-stage NLP pipeline: (i)~intent classification for agricultural commands;
(ii)~slot filling with domain-grounded entity resolution; (iii)~command compilation to
executable robotic API calls.

\textbf{Input data.}
A curated conversational corpus built by interviewing agronomists at partner farms (ProBayes, UniMI);
augmented with GPT-4 synthetic dialogues subject to human expert validation.

\textbf{Expected output.}
A fine-tuned LLM (e.g., Llama-3 Instruct, Mistral-7B) with PEFT adapters trained on
agricultural intent types: \texttt{QUERY}, \texttt{COMMAND}, \texttt{ABORT}, \texttt{REPORT},
\texttt{DIAGNOSE}. Output includes a confidence score and a structured command object
(JSON-serialisable) for the robotic layer.

\textbf{Relation to overall objective.}
This is the core NLP contribution. My existing expertise in LLM fine-tuning
(Llama-3.2, Gemma, Mistral, Qwen — see CV), mechanistic interpretability,
and conversational AI directly positions me to execute and publish this work at ACL/EMNLP.

% ---- Task 3 ---------------------------------------------------------
\subsection{Task 3 — Metaflow MLOps Pipeline: Prototype $\to$ Edge $\to$ Production
  \textnormal{[Months 6--18, UCBLyon1 + ProBayes secondment]}}

\textbf{Objective.}
Implement the full Metaflow-orchestrated ML pipeline following the prototype-to-production
lifecycle described in Fig.~1 of~\cite{tagliabue_metaflow_2023}: local prototyping $\to$
cloud training $\to$ edge deployment $\to$ scheduled production inference.

\textbf{Input data.}
Artifacts from Tasks 1 and 2; model checkpoints from fine-tuning runs.

\textbf{Expected output.}
A reproducible Metaflow \texttt{FlowSpec} DAG with: (a)~\texttt{@batch} decorator for GPU-based
LLM fine-tuning on cloud; (b)~\texttt{@card} auto-documentation; (c)~\texttt{@kubernetes}
edge deployment targeting Raspberry Pi 4 / NVIDIA Jetson clusters at field stations;
(d)~\texttt{@schedule} for daily model refresh from new field data.

\textbf{Relation to overall objective.}
This task directly implements the Metaflow architecture: Runtime, Data Store (S3),
Metadata Service, and Client API as described in~\cite{tagliabue_metaflow_2023}.
The ProBayes secondment (months 12--18) provides industry validation of the pipeline.

% ---- Task 4 ---------------------------------------------------------
\subsection{Task 4 — Robot Integration: Conversational Orchestration over Spraying Missions
  \textnormal{[Months 13--24, UniMI]}}

\textbf{Objective.}
Integrate the AgriTalk conversational layer with UniMI's real robotic platforms, enabling
natural language orchestration of pesticide spraying missions including start/stop,
re-routing, dosage adjustment, and failure diagnosis.

\textbf{Input data.}
Real robotic telemetry from UniMI's precision spraying robots; GPS field maps;
crop health indices from aerial multispectral cameras.

\textbf{Expected output.}
A ROS2-integrated middleware that: (a)~receives compiled commands from Task 2;
(b)~validates them against live robot state and field constraints;
(c)~requests HITL confirmation for safety-critical actions;
(d)~provides natural language explanations of robot decisions via the conversational interface.

\textbf{Relation to overall objective.}
This is the primary agricultural science contribution, grounded in Prof.\ Oberti's
precision crop protection research~\cite{oberti_2025} and the multi-modal sensing
advances discussed in~\cite{chen_2021}.

% ---- Task 5 ---------------------------------------------------------
\subsection{Task 5 — Responsible AI Audit, User Study \& Longitudinal Evaluation
  \textnormal{[Months 20--36, UniMI + UCBLyon1]}}

\textbf{Objective.}
Conduct a structured user study with real farmers and agronomists, measuring trust calibration,
over-reliance risks, explainability quality, and system usability under real field conditions.

\textbf{Input data.}
Deployment logs from Task 4; user interaction transcripts; expert annotations.

\textbf{Expected output.}
A published evaluation framework for conversational AI in safety-critical agricultural settings,
including a set of \textbf{Responsible AI metrics} (defined in Section 5) and a
\textbf{DAG Card} audit trail~\cite{tagliabue_dag_card_2021} for every spraying decision.

\textbf{Relation to overall objective.}
This closes the feedback loop: user corrections feed back into Task 2 fine-tuning
via the Metaflow pipeline (Task 3), creating a self-improving, human-in-the-loop system.

% =====================================================================
\section{3.\ Workflow and Infrastructure: Metaflow as the Research Backbone}
% =====================================================================

The Metaflow framework~\cite{tagliabue_metaflow_2023,tuulos_2022} is the central
infrastructure choice for AgriTalk. The rationale is precise: agricultural ML pipelines share exactly
the pathologies that Metaflow was designed to cure — heterogeneous compute needs (CPU edge
nodes for inference, GPU cloud for training), wide gap between experiment and production,
and multi-team ownership (NLP team at Lyon, robotics team at Milan).

\vspace{4pt}
\textbf{Architecture decisions:}
\begin{itemize}[leftmargin=1.2em,itemsep=2pt]
  \item \textbf{Data Store}: MinIO (local/edge) mirrored to AWS S3 (cloud). All IoT sensor
  artifacts, model weights, and command logs are content-addressed for reproducibility.
  \item \textbf{Compute}: \texttt{@local} for rapid prototyping; \texttt{@batch} (AWS) for
  LLM fine-tuning; \texttt{@kubernetes} on NVIDIA Jetson for edge inference.
  \item \textbf{Orchestration}: \texttt{@schedule} daily retraining from new field data;
  \texttt{AWS Step Functions} for production DAG execution.
  \item \textbf{Versioning}: Every Metaflow \texttt{Run} is tagged with \texttt{field\_id},
  \texttt{crop\_type}, \texttt{model\_version}, enabling precise reproducibility of any past decision.
  \item \textbf{Documentation}: \texttt{@card} generates shareable DAG cards~\cite{tagliabue_dag_card_2021}
  for every sprint, creating an audit trail that can be reviewed by supervisors and ethics boards.
  \item \textbf{Streaming integration}: Apache Kafka topics feed into Metaflow \texttt{@step}s
  via \texttt{foreach} branches — one branch per sensor modality — enabling parallel
  horizontal scaling of IoT stream processing.
\end{itemize}

\textbf{Backup plan.} Should Metaflow's Kubernetes integration encounter issues on UniMI's
on-premises infrastructure, the fallback is Apache Airflow 2.x with MLflow for artifact
versioning — a functionally equivalent architecture with higher DevOps overhead~\cite{tagliabue_metaflow_2023}.

% =====================================================================
\section{4.\ Methods and Experiments}
% =====================================================================

\subsection{4.1 NLP \& LLM Methods}

\textbf{Fine-tuning strategy.} I will apply Parameter-Efficient Fine-Tuning (PEFT/LoRA) to
Llama-3-8B-Instruct and Mistral-7B-Instruct on the agricultural command corpus (Task 2).
Given my existing hands-on experience with these exact model families and LoRA fine-tuning
(CV — BVF project), ramp-up time is minimal.

\textbf{Semantic grounding.} To prevent hallucination of pesticide names and dosages,
the LLM is constrained by a Retrieval-Augmented Generation (RAG) layer backed by
EPPO Plant Health and PPDB pesticide databases. Output tokens are scored for domain
consistency before compilation.

\textbf{Uncertainty quantification.} Inspired by my work on belief geometry in LLMs
(BVF project, CV), I will apply conformal prediction~\cite{shafer_vovk_2008} and
entropy-based confidence scoring to flag low-confidence commands for mandatory HITL review.
This is a direct bridge between my current mechanistic interpretability research and
the agricultural application domain.

\subsection{4.2 Systems and MLOps Methods}

\textbf{Streaming pipeline.} Apache Spark Structured Streaming processes IoT data at
$<$500ms latency. Field events (sensor anomalies, weather alerts) trigger Metaflow
\texttt{@trigger} re-evaluation of current spray missions.

\textbf{Edge inference.} Quantized (4-bit GGUF) models run on NVIDIA Jetson Orin at field
stations. Latency budget: $\leq$200ms for intent classification; full chain $\leq$800ms.

\textbf{Federated privacy.} Drawing on my federated learning research (CFL, published
SpringerNature ICOMP'24), farm-specific model adaptation will use federated fine-tuning
to avoid centralising sensitive crop yield and field geometry data.

\subsection{4.3 Datasets and Environments}

\begin{tabularx}{\textwidth}{@{}lXl@{}}
  \toprule
  \textbf{Dataset / Environment} & \textbf{Purpose} & \textbf{Phase} \\
  \midrule
  ACRE Challenge (drone imagery) & Crop health indexing & Tasks 1, 4 \\
  OpenFarm / PPDB / EPPO & RAG knowledge base & Task 2 \\
  FieldRobot Event & Robot kinematic simulation & Task 4 \\
  UniMI Spraying Robot Logs & Real-world integration test & Task 4 \\
  ROS2 Gazebo Agricultural Worlds & Safe pre-deployment simulation & Tasks 3, 4 \\
  Custom Agronomist Dialogue Corpus & Intent fine-tuning & Task 2 \\
  \bottomrule
\end{tabularx}

% =====================================================================
\section{5.\ Evaluation Criteria and Metrics}
% =====================================================================

\textbf{Technical metrics:}
\begin{itemize}[leftmargin=1.2em,itemsep=1pt]
  \item \textbf{Intent accuracy:} F1 score on 5-class command taxonomy; target $>$92\% on held-out agronomist dialogues.
  \item \textbf{Command safety rate:} \% of ABORT / safety-critical commands correctly intercepted before execution; target 100\%.
  \item \textbf{End-to-end latency:} NL input $\to$ robot actuation; target $\leq$800ms at edge.
  \item \textbf{Pipeline reproducibility:} All Metaflow runs must be deterministically reproducible
  from stored artifacts; validated by replaying 20\% of historical runs~\cite{tagliabue_metaflow_2023}.
  \item \textbf{Calibration error:} Expected Calibration Error (ECE) of confidence scores; target ECE $<$0.05.
  \item \textbf{Pesticide reduction:} In simulation, target $\geq$25\% reduction vs.\ uniform spray baseline.
\end{itemize}

\textbf{User-centred metrics:}
\begin{itemize}[leftmargin=1.2em,itemsep=1pt]
  \item \textbf{System Usability Scale (SUS):} Target $>$75 (``Good'') for non-expert farmers.
  \item \textbf{Trust calibration:} User perceived reliability vs.\ actual reliability (Pearson $\rho$); well-calibrated means $\rho > 0.8$~\cite{kaur_2020}.
  \item \textbf{Explanation quality:} Likert-scale rating by agronomists on decision justifications; target $>$4/5.
  \item \textbf{Over-reliance rate:} Frequency of users accepting AI recommendations without review when uncertainty is high; target $<$15\%.
\end{itemize}

% =====================================================================
\section{6.\ Responsible AI Safeguards}
% =====================================================================

\begin{mdframed}[backgroundcolor=red!5, linewidth=0.5pt]
\textbf{Core principle.} AgriTalk is an \emph{interface}, not an autonomous agent.
No robotic actuation occurs without either explicit human confirmation or a pre-approved
autonomous-safe policy validated by a certified agronomist.
\end{mdframed}

The six specific safeguards are:

\begin{enumerate}[leftmargin=1.2em,itemsep=3pt]
  \item \textbf{Mandatory HITL confirmation for irreversible actions.}
  Any command classified as \texttt{COMMAND} with action type \{\texttt{SPRAY}, \texttt{ABORT\_MISSION}, \texttt{DOSAGE\_CHANGE}\} triggers a structured confirmation dialogue before ROS2 command publication.
  Low-confidence commands (ECE-flagged) always route to HITL regardless of action type.

  \item \textbf{Uncertainty reporting and epistemic transparency.}
  Every system response includes a calibrated confidence bar (\textsc{High / Medium / Low})
  derived from conformal prediction bounds. ``I am not sure'' is a valid and encouraged
  system output~\cite{amershi_2019}.

  \item \textbf{Role-based permissions.}
  The system enforces three roles: \texttt{OBSERVER} (read-only queries), \texttt{OPERATOR}
  (read + low-risk commands), \texttt{AGRONOMIST} (full access). JWT-based authentication.
  Prompt injection attempts are sanitized via input validation before LLM processing.

  \item \textbf{Hallucination containment.}
  RAG grounding ensures pesticide names and dosages are sourced only from EPPO/PPDB
  databases. Any generated response that cannot be attributed to a retrieval source is
  flagged with a \textsc{[Unverified]} marker and blocked from actuation.

  \item \textbf{Audit trail and traceability.}
  Every NL query, compiled command, uncertainty score, human confirmation decision, and
  robot action is logged as a Metaflow versioned artifact with field ID, timestamp,
  and user ID. This enables post-hoc accountability review.

  \item \textbf{Fallback and degraded-mode operation.}
  If the LLM service is unavailable (edge network failure), the system falls back to
  a rule-based finite-state machine covering the 20 most common command templates.
  Farmers are clearly notified of the degraded mode.
\end{enumerate}

\textbf{Bias and fairness.} The agronomist dialogue corpus will be collected across
multiple EU farm types, crop varieties, and operator experience levels to prevent
bias toward large-scale intensive farming paradigms.

% =====================================================================
\section{7.\ Reflection on Tool Use}
% =====================================================================

I used GitHub Copilot (Claude Sonnet 4.6) as a \emph{structured thinking partner} in drafting
this proposal. Specifically: the tool helped me map the Metaflow architecture sections to
the paper's Fig.~1, Fig.~2, and Fig.~3; suggested the conformal prediction calibration
mechanism; and proposed the role-based permission taxonomy. I critically reviewed all outputs,
corrected several incorrect references (the tool initially hallucinated journal names), revised
the responsible AI section to better align with Prof.\ Oberti's precision protection
expertise~\cite{oberti_2025}, and rewrote the evaluation metrics entirely from my own experience
with LLM evaluation (BVF project). I take full ownership of every claim, citation, and
methodological choice in this proposal and am prepared to defend, revise, and critically discuss
all of it in the interview.

% =====================================================================
\section{8.\ Institutional Plan: 36-Month Roadmap}
% =====================================================================

\subsection{8.1 Year-by-Year Contribution Plan}

\textbf{Year 1 (UCBLyon1, Lyon, France) — Infrastructure \& NLP Core}

\noindent Months 1--4: Metaflow environment setup; IoT data ingestion DAG (Task 1);
literature review (MLOps + smart agriculture + conversational AI).\\
Months 5--8: Agricultural intent taxonomy design; corpus construction; AGROVOC ontology integration.\\
Months 9--12: LLM fine-tuning (Task 2); uncertainty quantification module.\\
\textbf{Target output:} 1 workshop paper (EMNLP-Workshop or ML4AG @ NeurIPS).

\textbf{Year 2 (UniMI, Milan, Italy + ProBayes Secondment) — Robotic Integration \& Production}

\noindent Months 13--18: ProBayes secondment — conversational prototype validation;
Metaflow production pipeline (Task 3); edge deployment on Jetson.\\
Months 19--24: ROS2 integration with UniMI robots; simulation experiments in Gazebo;
precision spraying mission orchestration (Task 4).\\
\textbf{Target output:} 1 full paper (ICRA / IROS) + 1 MLOps paper (MLSys or VLDB).

\textbf{Year 3 (UCBLyon1 + UniMI) — Evaluation, Responsible AI \& Thesis}

\noindent Months 25--30: User study with farmers; responsible AI audit; longitudinal evaluation (Task 5).\\
Months 31--36: Thesis writing; dissemination; one final paper (CHI / FAccT / ISWC).\\
\textbf{Target output:} PhD thesis; 1 Responsible AI paper; software artifact release (GitHub).

\subsection{8.2 Backup Plans}

\begin{tabularx}{\textwidth}{@{}lXX@{}}
  \toprule
  \textbf{Risk} & \textbf{Primary Plan} & \textbf{Backup Plan} \\
  \midrule
  LLM API unavailability & Cloud GPT-4 / Llama API & Local quantized model (GGUF) \\
  Robot access delay at UniMI & Gazebo + ROS2 simulation & FieldRobot Event virtual environment \\
  Metaflow Kubernetes issues & Metaflow on AWS ECS & Apache Airflow + MLflow \\
  Corpus data scarcity & Agronomist interviews & GPT-4 synthetic data + expert review \\
  Federated privacy constraints & Centralised w/ anonymisation & Differential privacy (OpenDP) \\
  \bottomrule
\end{tabularx}

\subsection{8.3 Publication Venues}

\begin{itemize}[leftmargin=1.2em,itemsep=1pt]
  \item \textbf{NLP/LLM:} ACL, EMNLP, NAACL, ACL-Findings, ML4AG Workshop @ NeurIPS
  \item \textbf{MLOps/Systems:} MLSys, VLDB, SysML, CIDR
  \item \textbf{Robotics:} ICRA, IROS, ECMR
  \item \textbf{Agriculture:} Computers and Electronics in Agriculture (Elsevier), Biosystems Engineering
  \item \textbf{Responsible AI:} FAccT, AIES, CHI, ISWC
\end{itemize}

% =====================================================================
\section{9.\ Why I Am the Right Candidate}
% =====================================================================

My background creates a rare and directly relevant intersection: LLM fine-tuning and
interpretability (BVF, nDNA projects), federated learning for privacy-sensitive data
(SpringerNature publication), teaching IoT at a government polytechnic (4 years, 50+ student
projects), industry-grade MLOps experience (IBM Watson cloud, IIT Kanpur threat intelligence
systems, Infosys bioinformatics pipelines), and active research in AI safety and alignment.

The GreenFieldData position is the convergence point for all of these threads:
it demands NLP expertise (my primary research), MLOps discipline (my industrial background),
IoT literacy (my teaching experience), and responsible AI commitment (my safety and alignment work).
Most critically, I bring the \emph{geometric} intuition about LLM internal states developed in the
BVF project — an intuition about \emph{when and why LLMs are uncertain} — that is
precisely what is needed to build calibrated, trustworthy conversational systems in safety-critical domains.

% =====================================================================
\bibliographystyle{ieeetr}
\bibliography{proposal_refs}
% =====================================================================

\end{document}
